\input{preamble-template.tex}
\hsslug{pd-trigger-lite-dip}
\begin{document}

\hstitleblock{Board guide \textperiodcentered\ September 2026}{pd-trigger-lite-dip}%
{A 25 by 21 mm USB-C PD trigger: flip one DIP switch to get 5, 9, 12, 15 or 20 V at up to 3 A on two solder pads.}

This board asks a USB-C Power Delivery charger for one fixed voltage and passes the charger's
VBUS straight to two 2.54 mm output holes. There is no regulator in the power path, so the output
is exactly what the charger delivers. You pick the voltage with a five-position DIP switch, one
switch per voltage, so changing it needs no soldering iron. It is the pd-trigger-lite board with
its solder links replaced by that switch, and 5 V added as a choice.

\hsplate{top.png}{Top view}{25 x 21 mm, 2 layers, 1 oz}{render}%
{USB-C receptacle on the left, the voltage switch SW1 at the top with the voltage printed beside
each position, the PD controller at the bottom, output pads on the right (+ is VBUS, - is ground).
The red LED sits under the output pads.}

\newpage
\section{Schematic}
\hslead{One controller, one capacitor, one switch and five resistors. Everything else is copper.}

The receptacle J1 is a 6-pin power-only part: two VBUS contacts, two ground contacts and the two
CC lines. The controller U1 (WCH CH224A) talks to the charger over CC1 and CC2. It is powered
from VBUS through the 0 ohm link R7, which only exists so the controller's pins can be reached
with thin copper, and C1 (1 uF, 50 V) decouples it. VBUS runs as a wide copper pour from J1 to
the output pad J2.

U1 reads its CFG1 pin at power-up. One side of every switch in SW1 is CFG1; the other side goes
to that position's resistor. Position 1 pulls CFG1 up to VHV through R5 (100k), which the CH224
manual lists as the way to request 5 V. Positions 2 to 5 tie CFG1 to ground through R1 to R4
(6.8k, 24k, 56k, 120k), the manual's single-resistor settings for 9, 12, 15 and 20 V. CFG2 and
CFG3 are left open; the chip pulls them up inside. R6 and D1 light the LED from VBUS.

\hsplate{schematic.png}{Schematic}{kicad/pd-trigger-lite-dip.kicad\_sch}{ERC 0 errors, 0 warnings}%
{Nets join by label name. SEL5 to SEL20 are the nodes between each resistor and its switch.}

\newpage
\section{Setting the voltage}
\hslead{Exactly one switch ON, set before you plug the charger in.}

\begin{hsblock}
\begin{tabularx}{\linewidth}{L{1.8cm}L{1.8cm}L{2.2cm}Y}
\hstoprule
\hshead{Switch ON} & \hshead{Output} & \hshead{Others} & \hshead{What it does} \\
\hstoprule
1 & 5 V  & all OFF & R5 100k pulls CFG1 up to VHV \\ \hline
2 & 9 V  & all OFF & R1 6.8k from CFG1 to ground \\ \hline
3 & 12 V & all OFF & R2 24k from CFG1 to ground \\ \hline
4 & 15 V & all OFF & R3 56k from CFG1 to ground \\ \hline
5 & 20 V & all OFF & R4 120k from CFG1 to ground \\ \hline
none & undefined & & CFG1 floats; do not use \\ \hline
two or more & not allowed & & can damage U1 (see below) \\
\hstoprule
\end{tabularx}
\end{hsblock}

On the board the voltage is printed beside each switch: 20, 15, 12, 9 and 5 from the top, so
position 1 (5 V) is the bottom one, nearest the output pads. A switch is ON when its slider is
pushed toward the side marked ON on the switch body, the edge opposite its printed 1 to 5. The
board's silkscreen marks the same side ``\textless ON'' (toward the voltage labels). Use a toothpick or a small screwdriver; the sliders are 1.27 mm apart.

\begin{hscallout}{Exactly one switch, set before plugging in}
The controller reads the switch once, when it powers up, so set it with the board unplugged and
re-plug after any change. With every switch OFF, CFG1 floats and the request is undefined. Never
turn on two at once: with switch 1 and any other ON, R5 and the ground resistor form a divider
from VHV, and once VHV rises above 5 V that can push CFG1 past its 3.8 V absolute maximum.
\end{hscallout}

\textbf{Output.} The two plated holes take 20 AWG wire or a 1x2 2.54 mm header, which you solder
yourself; JLC does not fit them. The rating is 3 A at any voltage. A charger without an e-marked
cable never offers more than 3 A, and the board makes no 5 A claim.

\textbf{LED.} D1 is lit whenever VBUS is present. Its current rises with the voltage, so it glows
faintly at 5 V and clearly at 12 V and above.

\textbf{First power-up.} Plug in with nothing on the output, measure the voltage across the pads,
and only then connect the load. The charger must offer the chosen voltage as a fixed PDO; many
phone chargers stop at 9 V, and a charger that lacks it leaves the output at 5 V. The board has no
fuse and no TVS diode, because the controller's datasheet asks for neither. A hot-plug spike at
20 V is the one plausible way to damage it, so at 20 V plug the charger into the wall first and the
cable into the board last.

\newpage
\section{What is on it}
\hslead{Twelve parts placed by JLC, three of them Extended. Everything else is bare copper.}

\begin{hsblock}
\begin{tabularx}{\linewidth}{L{1.9cm}Y L{2.0cm}L{1.6cm}R{0.8cm}}
\hstoprule
\hshead{Ref} & \hshead{Part} & \hshead{LCSC} & \hshead{Type} & \hshead{Qty} \\
\hstoprule
U1 & CH224A PD sink controller, ESSOP-10 & C42459160 & Extended & 1 \\ \hline
J1 & USB-C receptacle 6P, 3 A 20 V (TYPE-C-6M-001) & C2798175 & Extended & 1 \\ \hline
SW1 & DIP switch, 5 positions, 1.27 mm SMD (SHOU HAN 1.27-5P TPPT) & C7421518 & Extended & 1 \\ \hline
C1 & 1 uF 50 V X5R 0603 & C15849 & Basic & 1 \\ \hline
R1 & 6.8k 1\% 0603 (9 V) & C23212 & Basic & 1 \\ \hline
R2 & 24k 1\% 0603 (12 V) & C23352 & Basic & 1 \\ \hline
R3 & 56k 1\% 0603 (15 V) & C23206 & Basic & 1 \\ \hline
R4 & 120k 1\% 0603 (20 V) & C25808 & Basic & 1 \\ \hline
R5 & 100k 1\% 0603 (5 V) & C25803 & Basic & 1 \\ \hline
R6 & 10k 1\% 0603 & C25804 & Basic & 1 \\ \hline
R7 & 0 ohm 0603 & C21189 & Basic & 1 \\ \hline
D1 & red LED 0603 & C2286 & Basic & 1 \\
\hstoprule
\end{tabularx}
\end{hsblock}

Not assembled: J2, the pair of output holes. The three Extended parts each add JLC's per-part
feeder fee; no Basic part can replace any of them.

\section{What it costs}
\hslead{About 33 USD for five assembled boards before shipping, as an estimate.}

\begin{hsstatrow}[3]
\hsstat{27.42}{USD, JLC estimate for 5 boards: PCB 2.00, setup 8.00, 3 Extended fees 9.00, stencil 8.00, joints 0.42.}
\hsstat{5.44}{USD, parts for 5 boards at LCSC unit prices, about 1.09 per board (the switch is 0.53).}
\hsstat{32.9}{USD total for 5 boards, before shipping and tax.}
\end{hsstatrow}

These are estimates from the skill's price table (verified 2026-07-24), not a quote. Live JLC
prices have run 1.9 to 3.1 times higher than that table for bare PCBs, so expect the PCB line to be
nearer 4 to 6 USD. Shipping depends on the country and the carrier and is not included. The cart
at \hsurl{https://cart.jlcpcb.com/quote} is the only real price.

\hsprovenance{Estimate from fab/quote.json (order\_quote.py, qty 5, HASL, economic assembly) and
kicad/parts.json prices, read 2026-09-24.}

\clearpage
\section{Ordering at JLCPCB}
\hslead{Three files go up: the gerber zip, then BOM.csv and CPL.csv for assembly.}

\subsection*{Pre-buy list}
JLC's BOM review may show an Extended part as idle stock: unselected, at quantity 0, until it has
been bought into your JLC parts inventory. JLC's public parts search cannot tell idle stock from
stock that assembles straight away, so all three Extended parts are listed. For 5 boards:

\begin{hsblock}
\begin{tabularx}{\linewidth}{L{1.9cm}Y L{1.2cm}R{1.6cm}}
\hstoprule
\hshead{LCSC} & \hshead{Part} & \hshead{Ref} & \hshead{Qty to buy} \\
\hstoprule
C2798175 & USB-C receptacle 6P (TYPE-C-6M-001) & J1 & 5 \\ \hline
C7421518 & DIP switch 5P, 1.27 mm (1.27-5P TPPT) & SW1 & 5 \\ \hline
C42459160 & CH224A PD sink controller & U1 & 5 \\
\hstoprule
\end{tabularx}
\end{hsblock}

\hsprovenance{From fab/prebuy.csv (bom\_cpl.py, build of 5 boards).}

\begin{enumerate}
\item Open the JLC quote page (link above) and upload the gerber zip,
\texttt{fab/pd-trigger-lite-dip\_gerbers.zip}.
\item Check that the viewer reads 2 layers, 25 x 21 mm. Keep 1.6 mm, 1 oz, HASL, green, quantity 5.
\item Turn on PCB Assembly: Economic, top side, 5 boards.
\item Upload \texttt{fab/BOM.csv} and \texttt{fab/CPL.csv}. Confirm all twelve parts matched their
LCSC numbers and that J2 is absent, since it is not assembled.
\item If a part on the pre-buy list shows as unselected at quantity 0, buy it into your parts
inventory from the BOM review page, then select it and re-check the quantity.
\item In the placement preview, check four parts. SW1: its pin-1 mark sits at the bottom-right
corner, next to R5, so switch 1 is the 5 V row and the switch's ON edge faces the voltage labels. D1: its cathode mark must face the right board edge. U1: the pin-1 dot sits at
the corner nearest the capacitor C1. J1: the receptacle mouth points off the left edge. No
rotation correction was applied to the CPL, because every footprint comes from JLC's own library.
\item Run JLCDFM on the gerbers as a second opinion, then pay. Payment is yours to do.
\end{enumerate}

\hsprovenance{Gates on 2026-09-24: ERC 0/0, DRC 0 errors 0 unconnected, place, verify and DFM pass.}
\section{What it cost to design}
\hslead{About 6 USD of model spend went into designing this board. It is not part of the board's price.}

\begin{hsblock}
\begin{tabularx}{\linewidth}{Y R{2.2cm}}
\hstoprule
\hshead{Step} & \hshead{USD} \\
\hstoprule
Before the workspace (brief and skill read) & 1.13 \\ \hline
Schematic & 0.85 \\ \hline
Placement & 2.10 \\ \hline
DFM, order files and guide & 2.23 \\
\hstoprule
\textbf{Total} & \textbf{6.31} \\
\hstoprule
\end{tabularx}
\end{hsblock}

The split follows when hwde recorded each step, and the rows are rounded to the cent. A later round that wrote the pre-buy list and fixed this board's BOM cost 4.93 USD more, including an iteration that timed out. It is not in the total because it also built the pre-buy list feature for every board. The review-fix round resumed the design session, so its cost is already inside the total. Later skill-wide rebuilds of this guide are skill work and are not counted.

\hsprovenance{reports/cost.json (gen\_cost.py), priced from reference/model\_prices.yaml, which was fitted to Claude Code's own session totals on 2026-09-24.}
\end{document}
